%=========================================================================
%  FEEDBACK NACH DEM GESPRÄCH — optional.
%
%  Ein kurzes Dokument, das man nach einem Erstgespräch nachreichen kann.
%  Nutzt denselben Briefkopf und dasselbe Theme wie das Dossier.
%     make feedback
%=========================================================================

\documentclass[10pt, a4paper]{article}

\providecommand{\cvtheme}{kompakt}

\newcommand{\application}{2026-01-muster-ag}
\ifdefined\appoverride\let\application\appoverride\fi
\input{applications/\application/meta}
\providecommand{\cvjobtitle}{Entwicklungsingenieur}
\providecommand{\cvcontent}{cv-content}

\input{templates/themes/\cvtheme}
\input{templates/themes/\cvcontent}

\setlength{\parskip}{3pt}

\begin{document}

\makeletterhead

\noindent
{\color{headercolor}\sffamily\bfseries\large Feedback Erstgespräch -- \companyname, 1.~Januar 2026}

\lettersection{Gesamteindruck}

\noindent
Lorem ipsum dolor sit amet, consetetur sadipscing elitr, sed diam nonumy eirmod tempor invidunt ut labore et dolore magna aliquyam erat, sed diam voluptua.

\lettersection{Was mich überzeugt hat}

\noindent
At vero eos et accusam et justo duo dolores et ea rebum. Stet clita kasd gubergren, no sea takimata sanctus est Lorem ipsum dolor sit amet.

\lettersection{Was ich einbringe}

\noindent
Duis autem vel eum iriure dolor in hendrerit in vulputate velit esse molestie consequat, vel illum dolore eu feugiat nulla facilisis.

\lettersection{Offene Punkte für das nächste Gespräch}

\noindent
\textbf{Thema 1:} Lorem ipsum dolor sit amet, consetetur sadipscing elitr.

\textbf{Thema 2:} At vero eos et accusam et justo duo dolores et ea rebum.

\end{document}
